% Pillar-2 (ZVF) discussion + conclusion
\section{Discussion and Limitations}
\label{sec:p2-discussion}
ZVF earns its place as a descriptive diagnostic of signal starvation, but our
evidence argues against treating it as a standalone causal or incrementally
predictive statistic. (Our companion Pillar-7 takes up exactly this question as an
\emph{interventional} sequel---testing, rather than assuming, whether ZVF-derived
quantities can be made predictive and used for control; the two papers are
complementary, descriptive characterization here and falsifiable intervention there,
not contradictory.) Three limitations are load-bearing. First, ZVF is
mechanically coupled to reward sparsity, group size, and baseline accuracy, so its
correlations are partly bookkeeping rather than discovery; multivariate tests
against reward mean, entropy, and divergence proxies are needed before any causal
reading. Second, its strong association is with catastrophic \emph{collapse}, an
easy target, and only weak with the graded held-out outcome that practitioners
actually care about. Third, as an unsigned existence statistic it aliases mastery
($p \to 1$) with incapacity ($p \to 0$) and, as the cross-examination shows, can be
driven to zero by sub-reward jitter (for example a small length penalty), falsely
reporting a healthy batch. These are reasons to prefer magnitude- and sign-aware
diagnostics, which we outline but do not yet validate at the $\rho \geq 0.45$ bar
the cross-examination sets.

\section{Conclusion}
\label{sec:p2-conclusion}
The Zero-Variance Fraction is a useful, cheap, mechanistically grounded lens on
when binary-reward GRPO stops producing within-group signal, and it reliably
tracks collapse and drifts upward over training. It is not, on our data, a
standalone predictor of final generalization, and its unsigned form conflates two
opposite regimes. The constructive path forward is to measure the \emph{magnitude}
and \emph{sign} of within-group contrast rather than merely its existence, and to
test such replacements as early-window leading indicators against held-out outcome.
We release the per-group reward tensors precisely so that these stricter tests can
be run.
